\documentclass[slovak, twoside]{book}
\usepackage{scripture}
\usepackage[a5paper, outer=20mm, inner=20mm]{geometry}
\usepackage{graphicx}
\usepackage[program=lilypond]{lyluatex}
\usepackage[autocompile]{gregoriotex}
\usepackage{titlesec}
\usepackage{fontspec}
\usepackage{enumitem}
\usepackage{xparse}
\usepackage{tikz}
\usepackage{fancyhdr}
\usepackage[slovak]{babel}
\usepackage{lettrine}
\usepackage{etoolbox}
\usepackage{marginalia}

\marginaliasetup{
  pos=left, 
  width=10mm,
  xsep=0pt,
  yshift=0pt
}

\makeatletter
\patchcmd{\chapter}{\if@openright\cleardoublepage\else\clearpage\fi}{\clearpage}{}{}
\makeatother

\newcommand{\blankpage}{%
    \null
    \thispagestyle{empty}%
    \newpage
    }

\newcommand{\markup}[1]{%
    \indent\large\textcolor{red}{#1}%
    }

\newcommand{\psalmheading}[1]{%
    \begin{center}
    \noindent\large\textcolor{red}{#1}
    \end{center}
    }

\newcommand{\partheading}[1]{%
    \noindent\large\textcolor{red}{\textmd{\uppercase{#1}}}
}

\newcommand*{\annotation}[1]{
    \noindent\small\textcolor{red}{#1}\large
    \vspace{0.5em}
}

\newcommand*{\annotationcenter}[1]{
    \noindent\large\begin{center}\textcolor{red}{#1}\end{center}\large
}

\newcommand*{\predsedajuci}[1]{
    \noindent\textbf{#1}
    \large
}

\newcommand*{\euch}[1]{
    \noindent\Large\textbf{#1}\large
}

\newcommand*{\startletter}[1]{
    \noindent\Large\textbf{\textcolor{red}{#1}}\Large
}

% usage: \uarc[sep][depth]{text}
\NewDocumentCommand\uarc{O{.2ex} O{.3ex} m}{%
    \tikz[baseline=(arced node.base)]{
        \node [inner sep=0pt, outer sep=0pt] (arced node) {#3};
        \draw [transform canvas={yshift=-#1}] (arced node.south west) parabola [parabola height=-#2] (arced node.south east);
    }%
}

\NewDocumentCommand\uup{O{.6ex} O{.5ex} m}{%
    \tikz[baseline=(node.base)]{
        \node [inner sep=0pt, outer sep=0pt, anchor=base] (node) {#3};
        \draw [transform canvas={yshift=-#1}] (node.base west) -- ([yshift=#2]node.base east);
    }%
}

% 3. Klesanie (udown)
\NewDocumentCommand\udown{O{.6ex} O{.5ex} m}{%
    \tikz[baseline=(node.base)]{
        \node [inner sep=0pt, outer sep=0pt, anchor=base] (node) {#3};
        \draw [transform canvas={yshift=-#1}] ([yshift=#2]node.base west) -- (node.base east);
    }%
}

% 3. Episéma (hrubá vodorovná čiarka)
\NewDocumentCommand\uepi{O{.2ex} O{1.5pt} m}{%
    \tikz[baseline=(epi node.base)]{
        \node [inner sep=0pt, outer sep=0pt] (epi node) {#3};
        \draw [line width=#2, transform canvas={yshift=-#1}] (epi node.south west) -- (epi node.south east);
    }%
}

\scripturesetup{poetry/leftmargin=0pt}
\setlength\parindent{1em}
\setlength{\parskip}{0.5em} 
\gresetinitiallines{0}
\setmainfont{Minion Pro}

\titleformat{\section}
  {\color{red}\bfseries\MakeUppercase}{\thesection}{1em}{}

\titleformat{\chapter}{\bfseries\filcenter\LARGE}{}{0pt}{\MakeUppercase}
\titlespacing*{\chapter}{0pt}{0em}{0.5em}
\titlespacing*{\section}{0pt}{0em}{0.3em}
\assignpagestyle{\chapter}{fancy}

% Nastavenie štýlu strán
\pagestyle{fancy}
\fancyhf{} % Vymaže aktuálne nastavenia hlavičky a pätičky

% Párne strany (Even) -> vľavo (Left)
% [Číslo] [medzera 2em] [Názov kapitoly]
\fancyhead[LE]{\bfseries\thepage \hspace{2em} \bfseries\textcolor{red}{Omša svätenia olejov}}

% Nepárne strany (Odd) -> vpravo (Right)
% [Názov kapitoly] [medzera 2em] [Číslo]
\renewcommand{\chaptermark}[1]{\markboth{#1}{}}
\fancyhead[RO]{\bfseries\textcolor{red}{\leftmark} \hspace{1em} \thepage}
\setlength{\headheight}{24pt}

% Nastavenie čiary
\renewcommand{\headrulewidth}{0.4pt} % Hrúbka čiary pod hlavičkou
\renewcommand{\footrulewidth}{0pt}   % Bez čiary nad pätičkou


\begin{document}
    \begin{titlepage}
        \begin{center}
            \begin{huge}\textcolor{red}{\textbf{MISSA CHRISMATIS}}\end{huge}
            \\\textbf{OMŠA SVÄTENIA OLEJOV}
            \vspace*{0.5cm}
            \\ \includegraphics[width=9cm]{./missa_chrismatis.png}
            \vspace*{0.5cm}
            \large
            \\ \textbf{Katedrála svätého Martina}
            \\ \textit{Spišská Kapitula - Spišské Podhradie}
        \end{center}

    \end{titlepage}

\newpage
\blankpage
\large
\textbf{Texty v tejto brožúre, prevzaté z oficiálnych liturgických kníh, sa zhodujú s pôvodinou 
a sú vydané pre potrebu Katedrály sv. Martina v Spišskej Kapitule — Spišskom Podhradí.} 
\vspace*{\fill}
\begin{center}
    \textbf{Kňazský seminár biskupa Jána Vojtaššáka}
    \\ \textbf{Spišská Kapitula — Spišské Podhradie}
    \\ \textbf{MMXXVI}
\end{center}
\newpage
\chapter{Úvodné obrady a liturgia slova}
\noindent\textcolor{red}{\small (FANFÁRY)}
\section*{Introit}
\normalsize
\lilypondfile{./Noty/vstupny_spev.ly}
\large
\begin{enumerate}[label=\color{red}\bfseries\theenumi.]
    %Treba dolpniť podčiarknutia a zvýraznenia pre spev a tie kvačky pod slabikami...
    \item Mo\textbf{je} srdce prekypuje krásnymi slovami, \textcolor{red}{\textbf{\GreDagger}}\space svoje verše ve\textbf{nu}jem \underline{krá}ľo\textbf{vi.} \textbf{\textcolor{red}{\gresixstar}} Môj \textbf{ja}zyk je ako pero \textbf{rých}lo\uarc{pis}ca. \textbf{\textcolor{red}{— \Rbar}}
    \item Ty \textbf{si} najkrajší z ľudských synov. \textbf{\textcolor{red}{\GreDagger}}\space Z tvojich perí \textbf{ply}nie underline{mi}lo\textbf{ta.} \textbf{\textcolor{red}{\gresixstar}} Pre\textbf{to} ťa Boh po\textbf{žeh}nal \underline{na}veky. \textbf{\textcolor{red}{— \Rbar}}
    \item Ty \textbf{naj}mocnejší, pripáš \textbf{si} meč \underline{na} bed\textbf{rá;} \textbf{\textcolor{red}{\gresixstar}} svo\textbf{ju} vele\textbf{bu} a \underline{dôs}tojnosť. \textbf{\textcolor{red}{— \Rbar}}
    \item Vo \textbf{svo}jej dôstojnosti šťastne vytiahni, nasadni na voz \textbf{\textcolor{red}{\GreDagger}}\space a bojuj zapravdu, lásku a \textbf{spra}vod\underline{li}\textbf{vosť.} \textbf{\textcolor{red}{\gresixstar}} Nech \textbf{ťa} tvoja pravica učí konať ú\textbf{žas}né \uarc{skut}ky. \textbf{\textcolor{red}{— \Rbar}}
    \item Tvo\textbf{je} ostré šípy \textbf{\textcolor{red}{\GreDagger}}\space zasiahnu srdcia kráľových \textbf{ne}pria\underline{te}\textbf{ľov;} \textbf{\textcolor{red}{\gresixstar}} pod\textbf{da}jú \textbf{sa} ti \underline{ná}rody. \textbf{\textcolor{red}{— \Rbar}}
    \item Tvoj \textbf{trón,} Bože, \textbf{tr}vá \underline{na}ve\textbf{ky} \textbf{\textcolor{red}{\gresixstar}} a \textbf{žez}lo tvojho kráľovstva je žezlo \textbf{spra}\-vod\underline{li}vosti. \textbf{\textcolor{red}{— \Rbar}}
    \item Mi\textbf{lu}ješ spravodlivosť a nenávidíš neprávosť, \textbf{\textcolor{red}{\GreDagger}}\space preto ťa Boh, tvoj Boh, pomazal o\textbf{le}jom \underline{ra}do\textbf{sti} \textbf{\textcolor{red}{\gresixstar}} viac \textbf{a}ko \textbf{tvo}jich \uarc{dru}hov. \textbf{\textcolor{red}{— \Rbar}}
    \item Tvoj \textbf{o}dev vonia myrhou, alo\textbf{ou} a \underline{ka}si\textbf{ou} \textbf{\textcolor{red}{\gresixstar}} a \textbf{roz}veseľuje ťa zvuk harfy z palácov zo \textbf{slo}no\uarc{vi}ny. \textbf{\textcolor{red}{— \Rbar}}
\end{enumerate}

\noindent\textcolor{red}{\textbf{HYMNA JUBILEA}}
\vspace{0.5em}
{\normalsize \lilypondfile{./Noty/jub1.ly}}
\begin{enumerate}[label=\color{red}\bfseries\theenumi.]
    \item Sláva ti, Bože, za Cirkev živú, ktorá v láske, jednote stojí,
    Duch Svätý dvíha nás, posilňuje, v srdciach verných oheň tvoj horí.
    \item Hlas tvojho slova v našej zemi znie, vďaka svedkom viery stále žije v nás,
    vo svetle ich krokov verne kráčame, tvoje dielo rastie v každom z nás.
    \item Ako svätý Martin, pastier pokorný, každý teba nech zviditeľniť vie,
    plášťom milosrdenstva odievať mnohých, to je cesta, v ktorej Cirkev zrie.
    \item Túžime šíriť svetlo nádeje, v tmách sveta aj dnes lásku tvoju niesť,
    nech domom pokoja tvoja Cirkev je, znakom viery pre celý svet. 
\end{enumerate}

\noindent\textcolor{red}{\textbf{JKS 244}}
\begin{enumerate}[label=\color{red}\bfseries\theenumi.]
    \item \textbf{K stolu božej láskavosti} - s pokorou prichádzame, - pred Ježišom v skrúšenosti - na kolená padáme, - [:keď večere Pána slávnej, - lásky jeho k svetu dávnej, - dnes pamiatku v úcte zjavnej - slávne vykonávame.:]
    \item Tebe slávu bez skončenia, - Kriste, prespevujeme, - pre večeru tú spasenia - teba zvelebujeme. - [:Ty, nesmiernej velebnosti, - prijmi naše povďačnosti, - ktoré tvojej láskavosti - teraz prisluhujeme.:]
    \item Veríme ťa prítomného - vo Sviatosti Oltárnej, - Bohčloveka láskavého, - čo v pamiatke preslávnej - [:pod spôsobom vína, chleba - k požívaniu dávaš seba - nám, veriacim, Pane z neba, - z lásky k hriešnym nám zjavnej.:]
\end{enumerate}

\newpage
\section*{Kyrie}
\hfill{\small LS I 665}\par
{\normalsize \lilypondfile{./Noty/kyrie.ly}}
\vspace{0.5em}
\section*{Oslavná pieseň - glória}
\hfill{\small LS I 658}\par
{\normalsize \lilypondfile{./Noty/gloria.ly}}

\newpage
\section*{Medzispev}
\hfill{\small Katarina vl.}\par
{\normalsize \lilypondfile{./Noty/zalm.ly}}

\hfill{\small gregoriánsky I. tónus}\par
{\normalsize \lilypondfile{./Noty/aleluja.ly}}

\chapter{Obnovenie kňazských sľubov}
\annotation{Odpovedajú len prítomní kňazi.}
\annotation{Po homílii sa biskup prihovorí takto alebo podobne:}

\predsedajuci{Milovaní bratia, pri výročnej spomienke na deň, keď Kristus
Pán dal apoštolom i nám účasť na svojom kňazstve, sa vás
pýtam: Chcete pred svojim biskupom a pred Božím ľudom
obnoviť sľuby, ktoré ste kedysi urobili?}

\annotation{Všetci kňazi naraz odpovedia:}
\\Chcem.

\predsedajuci{Chcete sa bližšie primknúť k Ježišovi, nášmu Pánovi, a lepšie
sa mu prispôsobiť, zrieknuť sa samých seba a verne plniť
kňazské poslanie, ktoré ste z lásky ku Kristovi a jeho Cirkvi
vzali na seba v deň svojej kňazskej vysviacky?}

\annotation{Všetci kňazi:}
\\Chcem.

\predsedajuci{Chcete v Eucharistii i v ostatných liturgických úkonoch
verne vysluhovať Božie tajomstvá a vykonávať posvätný
učiteľský úrad podľa príkladu Ježiša Krista, Hlavy a Pastie-
ra, teda nie z túžby po hmotných výhodách, ale jedine
z lásky k dušiam?}

\annotation{Všetci kňazi:}
\\Chcem.

\annotation{Potom sa biskup obráti k ľudu a pokračuje.}
\\\annotation{Odpovedajú spoločne všetci prítomní:}

\predsedajuci{A vy, milovaní bratia a sestry, modlite sa za svojich kňazov,
nech ich Pán zahrnie svojimi darmi, aby vás ako verní
služobníci Krista, najvyššieho kňaza, viedli k nemu,
k prameňu spásy.}
\annotation{Ľud:}
\\Kriste, uslyš nás. Kriste, vyslyš nás.

\predsedajuci{Modlite sa aj za mňa, aby som verne plnil svoje apoštolské
poslanie a bol medzi vami každý deň živým a dokonalejším
obrazom Krista kňaza, dobrého pastiera, učiteľa a služobníka
všetkých.}

\annotation{Ľud:}
\\Kriste, uslyš nás. Kriste, vyslyš nás.

\predsedajuci{Nech nás Pán zachová vo svojej láske a nech nás všetkých
kňazov i ľud, privedie do večného života.}

\annotation{Všetci:}
\\Amen.

\section*{Hymnus}
{\normalsize \lilypondfile{./Noty/vykupitel_prijmi.ly}}
\begin{enumerate}[label=\color{red}\bfseries\theenumi.]
    \item Plody stromu, k posväteniu živným slnkom zrodené, prináša ľud
    zbožný tomu, čo dal svetu spasenie. \textbf{\textcolor{red}{— \Rbar}}
    \item Vznešený kráľ večnej vlasti, posväť teda láskavo olej na znak, že si
    obral zlého ducha o právo. \textbf{\textcolor{red}{— \Rbar}}
    \item Novú silu daruj krizmou pomazaným vo viere, by zranená dôstojnosť
    sa skvela v novej nádhere. \textbf{\textcolor{red}{— \Rbar}}
    \item A keď krstná voda z duše všetky viny vyhostí, na čele kríž svätou
    krizmou vleje do nej milosti. \textbf{\textcolor{red}{— \Rbar}}
    \item V lone Otca narodený z lona Matky ľudský syn, krizmou daj lúč, uväzni
    smrť, pomstiteľku našich vín. \textbf{\textcolor{red}{— \Rbar}}
    \item Tento sviatok každoročne vnáša radosť medzi nás, nech je zasvätený
    chvále, ktorú neumlčí čas. \textbf{\textcolor{red}{— \Rbar}}
\end{enumerate}

\chapter{Liturgia eucharistie}

\section*{Offertórium}
Prijmi, Pane, obetu, ktorú Ti tu skladám,
naše malé snaženia, ktorých veľkosť hľadám.
Verím, Pane, že si tu, každý deň a všade,
všetky naše tajomstvá, nech sú v Tvojej sláve.

\noindent Všetkým nám, dal si dar, zákon svojej lásky.
Všetkých nás, máš nás rád, nech Ti spieva každý.

\section*{JKS 244/4}
\begin{enumerate}[label=\color{red}\bfseries\theenumi.]
    \setcounter{enumi}{3}
    \item Oltár je, hľa, tvoj stôl svätý, - kde ti obetuje kňaz - chleba, vína dve podstaty, - jak si nám to kázal raz; - [:obeť veľkej velebnosti - ruší dávne pamätnosti, - pred tvár tvojej láskavosti - slávnostne pristupuje.:]
    \item Naveky si, svätý Bože, - neba, zeme ty Pane, - všetko ťa len chváliť môže, - velebiť neprestajne. - [:V prevelebnej tu Sviatosti, - v dare toľkej láskavosti, - v svojej svätej prítomnosti - skrytý byť neprestaneš.:]
\end{enumerate}

\section*{Sanctus}
\hfill{\small LS I 629}\par
\normalsize
\lilypondfile{./Noty/sanctus.ly}

\newpage
{
    \Large
    \color{red}\bfseries\MakeUppercase{\centering Prvá eucharistická modlitba\par (Rímsky kánon)\par}
}

\annotation{Hlavný celebrant rozopne ruky a povie:}\\
\noindent\euch{\startletter{T}eba teda, najláskavejší Otče, \marginalia{\textcolor{red}{Hl}}
\\pokorne vzývame a prosíme
\\skrze Ježiša Krista,
\\tvojho Syna a nášho Pána:}

\annotation{Zopne ruky a povie:}
\\\euch{láskavo prijmi}

\annotation{Urobí jeden znak kríža spolu nad chlebom a kalichom, hovoriac:}
\\\euch{a požehnaj \textcolor{red}{\grecross} tieto dary, tieto žertvy,
\\tieto sväté a nepoškvrnené obety.}

\annotation{Pokračuje s rozopätými rukami:}
\\\euch{\startletter{P}rinášame ti ich najmä
\\za tvoju svätú katolícku Cirkev
\\v jednote s tvojím služobníkom,
\\naším pápežom M., s naším biskupom M.*
\\a so všetkými, ktorí vyznávajú pravú,
\\katolícku a apoštolskú vieru.
\\Daruj svojej Cirkvi pokoj,
\\chráň ju, zjednocuj a spravuj po celom svete.}

\annotationcenter{Spomienka na živých}
\hskip-0.2em\euch{\startletter{P}amätaj, Pane, \marginalia{\textcolor{red}{1k}}
\\ na svojich služobníkov a služobnice \textcolor{red}{M.} a \textcolor{red}{M.}}

\annotation{Zopne ruky a chvíľu sa modlí za tých, ktorých si chce osobitne pripomenúť. Potom
s rozopätými rukami pokračuje:}
\\\euch{Pamätaj i na všetkých tu zhromaždených.
\\Ty poznáš ich vieru
\\a vieš, že sú ti oddaní.
\\Za nich ti prinášame túto obetu chvály.
\\Aj oni sami ti ju obetujú za seba
\\i za všetkých svojich drahých,
\\za svoje vykúpenie,
\\za nádej na večnú spásu i za časné blaho
\\a predkladajú svoje prosby tebe,
\\Bohu večnému, živému a pravému.}

\annotationcenter{Spomienka na svätých}
\hskip-0.2em\euch{\startletter{V} spoločenstve s oslávenou Cirkvou\marginalia{\textcolor{red}{2k}}
\\s úctou si spomíname
\\najmä na preblahoslavenú Máriu, vždy Pannu,
\\Rodičku Boha a nášho Pána Ježiša Krista,
\\ale i na svätého Jozefa, jej ženícha,
\\a na tvojich svätých apoštolov a mučeníkov
\\Petra a Pavla, Ondreja,
\\\textcolor{red}{(}Jakuba, Jána, Tomáša, Jakuba, Filipa,
\\Bartolomeja, Matúša, Šimona a Tadeáša,
\\Lína, Kléta, Klimenta, Sixta,
\\Kornela, Cypriána, Vavrinca, Chryzogóna,
\\Jána a Pavla, Kozmu a Damiána\textcolor{red}{)}
\\i na všetkých tvojich svätých.
\\Pre ich zásluhy a na ich prosby
\\poskytni nám vždy a všade svoju pomoc a ochranu.
\\\textcolor{red}{(}Skrze Krista, nášho Pána. Amen.\textcolor{red}{)}}

\annotation{Hlavný celebrant s rozopätými rukami pokračuje:}
\\\euch{\startletter{P}ane, milostivo prijmi túto obetu,\marginalia{\textcolor{red}{Hl}}
\\ktorú ti predkladáme my, tvoji služobníci,
\\i celá tvoja rodina.
\\Spravuj naše dni vo svojom pokoji,
\\zachráň nás od večného zatratenia
\\a pripočítaj k zástupu svojich vyvolených.}

\annotation{Zopne ruky.}
\\\euch{\textcolor{red}{(}Skrze Krista, nášho Pána. Amen.\textcolor{red}{)}}

\annotation{Vystrie ruky nad obetné dary a povie (i všetci koncelebranti vystrú ruku a povedia):}
\\\euch{Prosíme ťa, Bože,\marginalia{\textcolor{red}{Vš}}
\\láskavo urob túto obetu
\\vo všetkom požehnanou, dokonalou,
\\prijatou, duchovnou a ľúbeznou,
\\aby sa nám stala Telom a Krvou
\\tvojho milovaného Syna,
\\nášho Pána Ježiša Krista.}

\annotation{Zopne ruky.}
\\\annotation{V nasledujúcich formulách treba predniesť Pánove slová jasne a zrozumiteľne,
ako si to vyžaduje ich povaha.}
\\\euch{\startletter{O}n večer pred svojím umučením}

\annotation{Hlavný celebrant vezme chlieb, drží ho trocha pozdvihnutý nad oltárom a pokračuje (koncelebranti s ním):}
\\\euch{vzal chlieb
\\do svojich svätých a ctihodných rúk,}
\\\annotation{Pozdvihne oči.}
\\\euch{pozdvihol oči k nebu,
\\k tebe, Bohu, svojmu všemohúcemu Otcovi,
\\vzdával ti vďaky a dobrorečil, lámal chlieb
\\a dával svojim učeníkom, hovoriac:}

\annotation{Trocha sa skloní. (koncelebranti Pánové slová vyslovujú s pravicou vystretou smerom k chlebu, 
ak je to vhodné)}
\\\euch{
    \centering
    VEZMITE A JEDZTE Z NEHO VŠETCI,\\
    \centering
    LEBO TOTO JE MOJE TELO,\\
    \centering
    KTORÉ SA OBETUJE ZA VÁS.\\
}

\annotation{Hlavný celebrant ukáže konsekrovanú hostiu ľudu,
znova ju položí na paténu a pokľaknutím adoruje}
\\\annotation{Koncelebranti pri pozdvihovaní pozrú na hostiu a potom sa pri pokľaknutí hlavného celebranta hlboko uklonia.}

\annotation{Potom pokračujú:}
\\\euch{Podobne po večeri}

\annotation{Hlavný celebrant kalich, drží ho trocha zdvihnutý nad oltárom a spolu s koncelebrantmi pokračuje:}
\\\euch{vzal do svojich svätých a ctihodných rúk
\\aj tento preslávny kalich,
\\znova ti vzdával vďaky, dobrorečil
\\a dal ho svojim učeníkom, hovoriac:}

\annotation{Trocha sa skloní. (koncelebranti Pánové slová vyslovujú s pravicou vystretou smerom ku kalichu, ak je to vhodné)}
\euch{ \begin{center}VEZMITE A PITE Z NEHO VŠETCI,
\\LEBO TOTO JE KALICH MOJEJ KRVI,
\\KTORÁ SA VYLIEVA ZA VÁS I ZA MNOHÝCH
\\NA ODPUSTENIE HRIECHOV.
\\JE TO KRV NOVEJ A VEČNEJ ZMLUVY.
\\TOTO ROBTE NA MOJU PAMIATKU. \end{center}}

\annotation{Ukáže kalich ľudu, znova ho položí na korporál a pokľaknutím adoruje.}
\\\annotation{Koncelebranti pri pozdvihovaní pozrú na kalich a potom sa pri pokľaknutí hlavného celebranta hlboko uklonia.}

\annotation{Potom hlavný celebrant povie:}
\\\euch{\startletter{H}ľa, tajomstvo viery.}\marginalia{\textcolor{red}{Hl}}
\\\annotation{Ľud pokračuje zvolaním:}
\\Smrť tvoju, Pane, zvestujeme
\\a tvoje zmŕtvychvstanie vyznávame, kým neprídeš v sláve.
\\\annotation{Alebo:}
\\\euch{\startletter{V}yznajme tajomstvo viery.}\marginalia{\textcolor{red}{Hl}}
\\Pane, keď jeme tento chlieb a pijeme z tohto kalicha,
\\zvestujeme tvoju smrť, kým neprídeš v sláve.
\\\annotation{Alebo:}
\\\euch{\startletter{V}eľké je tajomstvo viery.}\marginalia{\textcolor{red}{Hl}}
\\Spasiteľ sveta, zachráň nás,
\\veď ty si nás vykúpil svojím krížom a zmŕtvychvstaním.

\annotation{Potom všetci koncelebranti spolu s hlavným celebrantom s rozpätými rukami hovoria:}
\\\euch{\startletter{P}reto, Pane, my, tvoji služobníci,\marginalia{\textcolor{red}{Vš}}
\\aj tvoj svätý ľud, slávime pamiatku
\\požehnaného umučenia a zmŕtvychvstania
\\i slávneho nanebovstúpenia Ježiša Krista,
\\tvojho Syna a nášho Pána,
\\a prinášame tebe, vznešenému Bohu,
\\dary z tvojich darov, svätý chlieb večného života
\\a kalich večnej spásy
\\ako obetu čistú, obetu svätú,
\\obetu nepoškvrnenú.}

\euch{\startletter{Z} hliadni na ne vľúdnym a láskavým okom
\\a milostivo ich prijmi,
\\ako si milo prijal obetné dary
\\svojho spravodlivého služobníka Ábela,
\\žertvu nášho praotca Abraháma
\\i svätú a nepoškvrnenú obetu
\\tvojho veľkňaza Melchizedecha.}

\annotation{Sklonení so zopätými rukami pokračuju:}
\\\euch{\startletter{P}okorne ťa prosíme, všemohúci Bože,
\\prikáž svojmu svätému anjelovi
\\preniesť tieto dary na tvoj nebeský oltár,
\\pred tvár tvojej božskej velebnosti,
\\aby nás všetkých,
\\ktorí máme účasť na tejto oltárnej obete
\\a prijmeme presväté Telo a Krv tvojho Syna,}

\annotation{Vzpriamia sa a prežehnajú, pritom hovoria:}
\\\euch{naplnilo hojné nebeské požehnanie a milosť.}

\annotation{Zopnú ruky.}
\\\euch{\textcolor{red}{(}Skrze Krista, nášho Pána. Amen.\textcolor{red}{)}}

\annotation{Spomienka na zosnulých}\\
\annotation{
Príhovory „Pamätaj, Pane“ za zomrelých a „Aj nás, svojich hriešnych služobníkov“ je vhodné zveriť jednému alebo dvom
koncelebrantom a každý sám prednesie tieto modlitby s rozopätými rukami a zvýšeným hlasom.
}
\\\euch{\startletter{P}amätaj, Pane,\marginalia{\textcolor{red}{1k}}
\\i na svojich služobníkov
\\a služobnice \textcolor{red}{M.} a \textcolor{red}{M.},
\\ktorí nás predišli do večnosti
\\so znakom viery
\\a spia spánkom pokoja.}
\\\annotation{Zopne ruky a všetci sa chvíľu modlia za zosnulých, ktorých si chcú
osobitne pripomenúť.}

\annotation{Potom s rozopätými rukami pokračuje:}
\\\euch{Prosíme ťa, Pane, daj im a všetkým,
\\ktorí odpočívajú v Kristovi,
\\prebývať na mieste blaha, svetla a pokoja.}

\annotation{Zopnú ruky.}
\\\euch{\textcolor{red}{(}Skrze Krista, nášho Pána. Amen.\textcolor{red}{)}}

\annotation{Pri slovách „Aj nás, svojich hriešnych služobníkov“ všetci koncelebranti sa bijú v prsia.}
\\\euch{\startletter{A}j nás, svojich hriešnych služobníkov,}\marginalia{\textcolor{red}{2k}}
\\\annotation{S rozopätými rukami pokračuje:}
\\\euch{ktorí dúfame v tvoje prehojné milosrdenstvo,
\\priveď do spoločenstva
\\svojich svätých apoštolov a mučeníkov:
\\Jána, Štefana, Mateja, Barnabáša,
\\\textcolor{red}{(}Ignáca, Alexandra, Marcelína, Petra,
\\Felicity, Perpetuy, Agáty, Lucie,
\\Agnesy, Cecílie, Anastázie\textcolor{red}{)}
\\a všetkých tvojich svätých.
\\Prosíme ťa, prijmi nás do ich spoločenstva,
\\nie pre naše zásluhy,
\\ale pre tvoje veľké zľutovanie.}
\\\annotation{Zopne ruky.}
\\\euch{Skrze nášho Pána Ježiša Krista.}

\newpage
\section*{\hfill POŽEHNANIE OLEJA NA POMAZANIE CHORÝCH\hfill}
\annotation{Skôr než biskup povie \textcolor{black}{„Skrze neho, ty, Bože“} v Prvej eucharistickej
modlitbe alebo doxológiu v Tretej eucharistickej modlitbe, diakon alebo
kňaz, donesie nádobu s olejom na pomazanie chorých k oltáru a drží ju
pred biskupom.}
\\\annotation{Biskup sa pri požehnaní oleja na pomazanie chorých modlí túto
modlitbu:}
\large
\\\predsedajuci{Bože, Otče, ty potešuješ všetkých
\\a chceš, aby chorí našli v tvojom Synovi útechu a zdravie.
\\Pozri na našu vieru a vypočuj naše prosby:
\\zošli z neba Utešiteľa Ducha Svätého
\\na tento olej — dar zelenej olivy —,
\\ktorý si nám dal na posilnenie tela,
\\a daj, nech tvoje sväté \textcolor{red}{\grecross} požehnanie
\\chráni telo, dušu i ducha každého,
\\kto prijme sviatosť pomazania chorých;
\\nech ho zbaví všetkých bolestí,
\\všetkých neduhov a chorôb.
\\Nech je nám tento olej požehnaný
\\v mene nášho Pána Ježiša Krista,
\\\textcolor{red}{(}ktorý s tebou žije a kraľuje na veky vekov.}

\noindent\textcolor{red}{\Rbar.:} Amen\textcolor{red}{)}.

\annotation{Záver \textcolor{red}{„ktorý s tebou žije“} sa používa len vtedy, keď sa toto
požehnanie koná mimo eucharistickej modlitby.}
\\\annotation{Po požehnaní oleja sa nádoba s olejom na pomazanie chorých
položí na pôvodné miesto a v omši sa pokračuje ako zvyčajne, až do
prijímania vrátane. Po požehnaní oleja hovorí sám hlavný celebrant:}

\euch{\startletter{S}krze neho ty, Pane,\marginalia{\textcolor{red}{Hl}}
\\všetky tieto dary stále tvoríš,
\\posväcuješ, oživuješ, požehnávaš a nám dávaš.}


\annotation{Hlavný celebrant vezme paténu s hostiou a jeden z koncelebrantov,
ak niet diakona, kalich, pozdvihnú ich a hovoria:}

\gregorioscore{./Noty/peripsum.gabc}
\annotation{Všetci:}
\\Amen (6x).

\chapter{Obrad prijímania}
\section*{Lámanie chleba}
\hfill {\small LS I 756}\par
{\normalsize \lilypondfile{./Noty/agnus.ly}}
\vspace{0.5em}
\section*{Spev na prijímanie}
{\normalsize \gregorioscore{./Noty/ave_verum.gabc}}

\section*{JKS 267}
\large
\begin{enumerate}[label=\color{red}\bfseries\theenumi.]
    \item Hostiu vítajme, - tej manne česť dajme, - zástupov Pánovi - a kráľov Kráľovi. - Zdrav’ buď, ó, anjelský chlebe, - voláme vrúcne ku tebe, - [:predrahé telo a krv.:]
    \item Národ náš posilňuj, - vlasť našu ochraňuj, - ku blahu všetkých nás - zošli nám pravdy jas. - Chráň nás vždy od nepriateľov, - pokoja nežičiteľov - [:Ježiša telo a krv.:]
    \item Daruj nám milosti, - ó, Bože, v štedrosti, - ktorý si chránieval - a mannou kŕmieval - oddaný ľud svoj v pravý čas, - napájaj krvou svojou nás: - [:ó, Bože milostivý:]
    \item A v smrti hodine - obdar nás, prosíme, - vierou i dúfaním - a pravým pokáním, - by sme tu prijali teba, - po smrti prišli do neba, - [:naveky chváliť teba.:]
    \item Zletujte, anjeli, - a všetci veselí - s nami sa spojujte, - Presvätý! spievajte. - Zdrav’ buď tu, Kriste, v sviatosti - prítomný, Bože z výsosti, - [:v spôsobe chleba, vína.:]
\end{enumerate}

\section*{Ó, srdce kňaszké kristovo}
\normalsize
\lilypondfile{./Noty/osrdceknazske.ly}
\large
\begin{enumerate}[label=\color{red}\bfseries\theenumi.]
    \item Ó, Srdce kňazské Kristovo, daj svätých kňazov ľudu nášmu, nech
    vzplanie láskou ich slovo, nech stvorí v ľudstve dušu krásnu. Ó,
    Srdce Krista atď. 
    \item Ó, Srdce kňazské Kristovo, bôľ lásky tvojej sa
    nás týka, veď odmrštia ťa surovo, náš hriech tkvie v tebe ako dýka.
    Ó, Srdce Krista atď. 
    \item Ó, Srdce kňazské Kristovo, nad nami zmiluj
    sa hriešnymi, hľa, každý tebe dá slovo: Viac už viac nebudeme
    zlými. Ó, Srdce Krista atď.
\end{enumerate}
\section*{ANTIFÓNA NA JUBILEUM}
{\normalsize \lilypondfile{./Noty/jub2.ly}}
\begin{enumerate}[label=\color{red}\bfseries\theenumi.]
    \item \uarc{Boh} vstá\udown{va} a jeho \uup{ne}priate\udown{lia} sa \uepi{tra}tia, \textbf{\textcolor{red}{\GreDagger}}\space spred jeho tváre utekajú \uup{tí,} čo ho \udown{ne}ná\uepi{vi}dia. \textbf{\textcolor{red}{— \Rbar}}
    \item \uarc{A}le \udown{spra}vodliví sa môžu tešiť a \uup{ja}sať pred \udown{Bo}žou tvárou \textbf{\textcolor{red}{\GreDagger}}\space a \uup{v ra}dos\udown{ti} sa \uepi{ve}seliť. \textbf{\textcolor{red}{— \Rbar}}
    \item \uarc{Chváľ}te \udown{Pá}na \uup{v je}\uarc{ho} \uepi{svä}tyni, \textbf{\textcolor{red}{\GreDagger}}\space chváľte ho \uup{na} jeho vzne\udown{še}nej \uepi{ob}lohe. \textbf{\textcolor{red}{— \Rbar}}
    \item \uarc{Chváľ}te \udown{ho} \uup{za} jeho \udown{či}ny \uepi{mo}hutné, \textbf{\textcolor{red}{\GreDagger}}\space chváľte ho \uup{za} jeho nes\udown{mier}nu \uepi{ve}lebnosť. \textbf{\textcolor{red}{— \Rbar}}
    \item \uarc{Chváľ}te \udown{ho} \uup{zvu}\uarc{kom} \uepi{poľ}nice, \textbf{\textcolor{red}{\GreDagger}}\space chváľte ho \uup{har}\udown{fou} a \uepi{ci}tarou. \textbf{\textcolor{red}{— \Rbar}}
    \item \uarc{Chváľ}te \udown{ho} \uup{bub}\udown{nom} a \uepi{tan}com, \textbf{\textcolor{red}{\GreDagger}}\space chváľte ho \uup{lý}\udown{rou} a \uepi{fla}utou. \textbf{\textcolor{red}{— \Rbar}}
    \item \uarc{Chváľ}te \udown{ho} ľubozvučnými cimbalmi, \uup{chváľ}te ho jasa\udown{vý}mi \uepi{cim}balmi; \textbf{\textcolor{red}{\GreDagger}}\space všetko, čo dýcha, \uup{nech} \udown{chvá}li \uepi{Pá}na. \textbf{\textcolor{red}{— \Rbar}}
    \item \uarc{Slá}va \udown{Ot}cu i Synu \uup{i} \udown{Du}chu \uepi{Svä}tému. \textbf{\textcolor{red}{\GreDagger}}\space Ako bolo na počiatku, tak nech je i teraz i vždycky \uup{i} na veky \udown{ve}kov. \uepi{A}men. \textbf{\textcolor{red}{— \Rbar}}
\end{enumerate}
\section*{Po prijímaní}

\predsedajuci{Všemohúci Bože,
\\ty nás obnovuješ svojimi sviatosťami; \textcolor{red}{\textbf{\gresixstar}}
\\vrúcne ťa prosíme, \textcolor{red}{\textbf{—}}
\\dopraj, aby sme sa mohli stať
\\príjemnou vôňou Ježiša Krista.
\\Lebo on žije a kraľuje na veky vekov.}

\textcolor{red}{Všetci:} Amen

\chapter{Požehnanie oleja katechumenov a posvätenie krizmy}
\section*{\hfill POŽEHNANIE OLEJA KATECHUMENOV\hfill}
\large
\annotation{Po modlitbe po prijímaní zíde biskup s asistenciou k pripravenému
stolu uprostred presbytéria. Všetci ostatní ostávajú na mieste. Biskup
najprv požehná olej katechumenov.}
\annotation{Biskup tvárou k ľudu a s rozopätými rukami sa modlí túto
modlitbu:}
\large
\\\predsedajuci{Bože, opora a mocná ochrana svojho ľudu,
\\ty si stvoril olej ako symbol sily;
\\láskavo ho \textcolor{red}{\textbf{\grecross}} požehnaj
\\a katechumenom, ktorí ním budú pomazaní,
\\daj odvahu, božskú múdrosť a silu,
\\aby hlbšie chápali Kristovo evanjelium
\\a odhodlali sa horlivo žiť kresťanským životom;
\\prijmi ich za svoje deti
\\a daj im nový, radostný život vo svojej Cirkvi.
\\Skrze Krista, nášho Pána.}

\noindent \textcolor{red}{Všetci:} Amen.

\section*{\hfill POSVÄTENIE KRIZMY\hfill}
\annotation{Biskup vleje vonné látky do oleja, a tak pripraví krizmu, ak už
predtým nebola pripravená. Keď to urobil, prednesie výzvu na
modlitbu.}

\predsedajuci{Milí bratia a sestry,
\\prosme všemohúceho Boha Otca,
\\aby požehnal a posvätil túto krizmu
\\a vo všetkých, ktorí ňou budú pomazaní,
\\nech vnútorné pôsobí Božia milosť,
\\aby dosiahli spásu.}

\annotation{Biskup môže dýchnuť na krizmu. Potom sa s rozopätými rukami
pomodlí jednu z nasledujúcich posväcujúcich modlitieb:}

\predsedajuci{Bože, od teba pochádza všetok vzrast
\\a duchovný pokrok;
\\láskavo prijmi ďakovnú službu,
\\ktorú ti našimi ústami a s radosťou prináša Cirkev.
\\Ty si na počiatku prikázal zemi,
\\aby vydala ovocné stromy, medzi nimi aj olivy,
\\darkyne tohto oleja na prípravu svätej krizmy.
\\Už prorok Dávid videl účinok tvojich sviatostí
\\a ospieval olej, ktorým naša tvár zažiari radosťou;
\\keď potopa očistila svet od hriechu,
\\holubica priniesla olivovú ratolesť,
\\čím naznačila budúcu úlohu olivy
\\a oznámila, že sa svetu vracia pokoj.
\\To sa v jasných účinkoch zjavuje aj dnes:
\\krstná voda zmýva nám všetky hriechy a priestupky
\\a tento olej dáva našim tváram radosť a pôvab.
\\Aj Mojžišovi si dal príkaz,
\\aby svojho brata Árona, vodou obmytého,
\\pomazal olejom za kňaza.
\\Ešte väčšej cti sa dostalo tomuto úkonu,
\\keď tvoj Syn a náš Pán Ježiš Kristus
\\požiadal Jána Krstiteľa, aby ho pokrstil v Jordáne.
\\Vtedy si z neba poslal Ducha Svätého v podobe holubice
\\a dosvedčil si,
\\že v tomto tvojom jedinom Synovi máš najväčšiu záľubu.
\\Ty si zjavne potvrdil Dávidovo proroctvo,
\\že Mesiáš bude pomazaný olejom radosti
\\viac ako ostatní ľudia.}

\annotation{Všetci koncelebranti vystrú pravú ruku ku krizme a mlčky ju držia
vystretú až do konca modlitby:}
\\\predsedajuci{Preto ťa, Pane, prosíme,
\\požehnaj a \textcolor{red}{\grecross} posväť tento olej
\\a vlož doň silu Ducha Svätého
\\a moc Ježiša Krista,
\\od ktorého krizma dostala meno.
\\Ako si pomazával kňazov, kráľov, prorokov,
\\tvojich svedkov,
\\tak nech táto krizma vo sviatostnom pomazaní
\\zavŕši dielo spásy
\\a dá plnosť života tým,
\\čo vychádzajú obnovení z krstného kúpeľa.
\\Posväť ich pomazaním,
\\zbav ich následkov prvotného hriechu,
\\aby sa stali chrámom tvojej velebnosti
\\a žiarili čistým a tebe milým životom.
\\A keď si ich touto sviatosťou
\\vyznačil kráľovskou, kňazskou a prorockou hodnosťou,
\\obleč ich do rúcha mravnej bezúhonnosti
\\a horlivosti v tvojej službe.
\\Nech je svätá krizma posilou pre spásu tým,
\\ktorí sa znovuzrodili z vody a z Ducha Svätého,
\\aby mali účasť na večnom živote v nebeskej sláve.
\\Skrze Krista, nášho Pána.}

\noindent\textcolor{red}{Všetci:} Amen.

\annotation{Po záverečnom požehnaní na konci omše vloží biskup tymian do
kadidelnice. Chór a ľud zatiaľ spieva pápežskú hymnu \textcolor{black}{\textit{V sedembrežnom kruhu
Ríma}}. Po skončení spevu asistujúci odídu v sprievode do sakristie. Posvätné
oleje nesú určení posluhujúci hneď za krížom.}

\section*{Na záver: JKS 523}
\begin{enumerate}[label=\color{red}\bfseries\theenumi.]
    \item V sedmobrežnom kruhu Ríma, - kde sa Petra chrám vypína. - Z tisíc
    hrdiel sa ozýva - pieseň nábožná, horlivá: - [:Živ, Bože, Otca Svätého, -
    námestníka Kristovho!:]
\end{enumerate}

\noindent\textcolor{red}{\small (FANFÁRY)}

\noindent\textcolor{red}{\textbf{JKS 394}}
\begin{enumerate}[label=\color{red}\bfseries\theenumi.]
    \item Ó, Mária bolestivá, - naša ochrana, - slovenský náš národ volá: - Pros
    za nás Boha. - [:Ty si Mať dobrotivá, - Patrónka ľútostivá, - oroduj vždy za
    náš národ - u svojho Syna.:]
    \item Matka drahá, spomni na nás - tu pod Tatrami, - verní sme my svätej
    Cirkvi, - prebývaj s nami! - [:Ty si Mať dobrotivá, - Patrónka ľútostivá, -
    oroduj vždy za náš národ - u svojho Syna.:]
    \item Po horách i po dolinách - venčíme tvoj chrám, - pútne miesta volajú
    ti: - Matka, národ chráň! - [:Ty si Mať dobrotivá, - Patrónka ľútostivá, -
    oroduj vždy za náš národ - u svojho Syna.:]
    \item Hlad, mor, vojny odstráň od nás, - vypros pokoja, - slovenský ľud
    nech vždy stráži - ochrana tvoja! - [:Ty si Mať dobrotivá, - Patrónka
    ľútostivá, - oroduj vždy za náš národ - u svojho Syna.:]
    \item Tebe sa tu oddávame, - Matka, pomôž nám, - by sme i my i potomci -
    vyhli vždy vojnám. - [:Ty si Mať dobrotivá, - Patrónka ľútostivá, - oroduj
    vždy za náš národ - u svojho Syna.:]
\end{enumerate}

\noindent\textcolor{red}{\textbf{JKS 503}}
\begin{enumerate}[label=\color{red}\bfseries\theenumi.]
    \item Lásky svätej ohnisko, - Srdce Boha Kráľa, - z oltára sa našich sŕdc - vznáša k tebe chvála! [:Božské Srdce, buď s nami, - kraľuj, panuj nad nami, - nech sa v našej rodine - život viery rozvinie.:]
    \item Božské Srdce, tvoji sme - v zasvätenia sľube, - zasvätení budeme - v tvojej žiť úľube. [:Božské Srdce, buď s nami, - kraľuj, panuj nad nami, - nech sa v našej rodine - život viery rozvinie.:]
    \item Z trónu tvojej dobroty - v náš dom zlietla láska, - v našom srdci vládnuť chce, - láskou milosť božská. [:Božské Srdce, buď s nami, - kraľuj, panuj nad nami, - nech sa v našej rodine - život viery rozvinie.:]
    \item V pokušení ochrana, - Srdce večnej sily. - Posilni nás, aby sme - v boji zvíťazili! [:Božské Srdce, buď s nami, - kraľuj, panuj nad nami, - nech sa v našej rodine - život viery rozvinie.:]
    \item Útočište hriešnikov, - Srdce zmilovania, - k tebe duša túli sa - v bôľny deň skonania. [:Božské Srdce, buď s nami, - kraľuj, panuj nad nami, - nech sa v našej rodine - život viery rozvinie.:]
\end{enumerate}

\newpage
{
\noindent
\hyphenpenalty=10000
\exhyphenpenalty=10000
Liturgické texty \textit{Omše svätenia olejov} boli doslovne prevzaté 
z tretieho vydania \textit{Rímskeho misála} v slovenskom jazyku (KBS, Bratislava 2021), 
\textit{Rímskeho pontifikála} v slovenskom jazyku (SSV Trnava 1981),
tretieho vydania \textit{Lekcionára I} v slovenskom jazyku (SSV Trnava 2003) 
a ďalších liturgických kníh a spevníkov na Slovensku.
}
\vfill
\noindent Zostavil:\\
Prof. ThDr. PhDr. Amantius Akimjak, PhD.,\\
člen LK KBS a LK Spišskej diecézy
\vfill
\noindent Recenzoval:\\
SLDr. Lukáš Fejerčák, PhD.,\\
Biskupský ceromoniár, tajomník diecézneho biskupa
\vfill
\noindent Sadzba:\\
Marek Villim\\
Ing. Jozef Pazúrik
\vfill
\noindent Pre vnútornú potrebu\\
Katedrály sv. Martina v Spišskej Kapitule
\vfill
\noindent Vydala:\\
Liturgická komisia Spišskej diecézy\\
Schola Cantorum Kňazského seminára biskupa Jána Vojtaššáka
\vfill
\noindent Ústav sakrálneho umenia TI TF KU\\
v Spišskej Kapitule v Spišskom Podhradí 2025
\end{document}